\Needspace{7\baselineskip}
\section{Relational Existence: The Foundational Axiom}
\label{sec:01}
Section 1 begins with the simplest possible starting point:

\begin{quote}
\textbf{There is something rather than nothing.}
\end{quote}
At first glance, this appears almost trivial. The central claim of this paper is that this starting point, once stated in its precise relational form, is strong enough to force a large amount of structure. The paper does not begin by assuming a manifold, a metric, a field, a background clock, or an external observer. It begins by asking what must be true before any of those notions can even be discussed.

The purpose of this section is therefore not to introduce physics directly. It establishes the admission rules that any later physics must satisfy. At the most fundamental level, no later observer is allowed to stand outside the relational system with access to arbitrary abstractions. Any later observer must be treated as part of the relational system it reads. Before continuous measurement, coordinate systems, or calibrated units, the primitive operation available at this layer is distinguishability: to separate one admitted content from what it is not, and to count what has been admitted.

\Needspace{5\baselineskip}
\subsection{The Existence Axiom}
The most basic question is simply:

\begin{quote}
\textbf{Is there something, or is there nothing?}
\end{quote}
If there is truly nothing, then there is nothing to describe, no distinction to make, no theory to construct, and no standpoint from which even the denial could be stated. A physical theory can only concern what exists. This leaves a boundary and an entry:

\begin{itemize}
\item \(0\): nothing exists. This is the boundary case. It admits no description and no theory.
\item \(1\): something exists. At least one distinction can be made.
\end{itemize}
Only the second case can carry a theory. But "something exists" already contains a distinction. To say that something exists is to distinguish it: from nothing, and from what it is not. Existence never arrives as a bare object to which distinguishability is later added. It arrives, if it arrives at all, as distinguishability within a relation. The paper therefore takes as its single physical axiom:

\begin{quote}
\textbf{Axiom 1 (relational existence).} Existence is relational distinguishability.
\end{quote}
The first number of the theory is not a second assumption. One primitive act of distinguishability is one binary cut, and only the admitted side of that cut enters the theory. The first admitted content is therefore exactly one distinction, with count

\begin{equation*}
\begin{adjustbox}{max width=\linewidth,center}
$\displaystyle
I = 1 .
$
\end{adjustbox}
\end{equation*}
Here \(I\) denotes the first existence count. This is not yet a statement about particles, space, time, fields, or geometry. It is the minimal assertion that there is something rather than nothing, made precise as distinguishability. That precision is what gives the rest of the theory something to build from.

\Needspace{5\baselineskip}
\subsection{Consequences of the Existence Axiom}
Once the first count \(I=1\) is admitted, the unit itself has a special status. It is not composed of smaller units of the same kind inside the theory, because then those smaller units would be the true primitives. The first unit of existence is therefore indivisible within this admission layer, and \(1\) is the natural primitive count.

Larger internal quantities arise only by counting further admitted units:

\begin{equation*}
\begin{adjustbox}{max width=\linewidth,center}
$\displaystyle
1,\ 2,\ 3,\ \dots
$
\end{adjustbox}
\end{equation*}
This is an ontological statement before it is a statement about numerical technique. The theory begins with admitted positive counts. Fractions, decimals, limits, and other numerical representations are not introduced here as new primitive existents. They may enter later as comparison reports, scale choices, limiting descriptions, or licensed mathematical language, but they are not the starting ontology of the theory.

This restriction concerns primitive admission, not the symbols available after a quantity has been derived. For an already established quantity, number class alone neither grants nor denies admissibility. An exact rational, irrational, algebraic, or transcendental expression remains a downstream representation of its typed object; it does not add another primitive relation-token. Its scientific warrant comes from the owning derivation, domain, and guards, not from the notation by itself.

Zero has a different status. It does not represent an existing object or quantity. It marks absence. Since the theory begins from the admitted case that there is something rather than nothing, zero is never treated as a primitive element of the ontology. It may be useful as an external boundary label or inside a controlled formal statement, but it is not an admitted object, primitive internal value, or held quantity. The bookkeeping keeps this distinction sharp: if no transfer occurred, the theory records that no transfer occurred, not a transfer of size zero; if a register is absent, it records absence, not a hidden holding of zero. A detector that did not fire and a detector that fired and read zero are different physical situations. Internally, presence is positive by type:

\begin{equation*}
\begin{adjustbox}{max width=\linewidth,center}
$\displaystyle
q \in \mathcal{D}_+ .
$
\end{adjustbox}
\end{equation*}
The same discipline applies to negative quantities. A negative sign is not interpreted as the existence of negative objects. It is a relational or oriented bookkeeping descriptor applied to positive content. If Alice has three apples and Bob receives two of them, nothing negative has been created. The relationship has changed: Alice now has one apple, Bob has two, and the transfer has an orientation. The internal formalism represents such orientation without turning negative objects into ontology.

A natural worry is that forbidding primitive negative quantities cripples the algebra physicists rely on, especially cancellation. It does not. Within a fixed kind of quantity, cancellation is not assumed as an additional physical axiom; it is a theorem:

\begin{equation*}
\begin{adjustbox}{max width=\linewidth,center}
$\displaystyle
a+q=b+q \Longrightarrow a=b .
$
\end{adjustbox}
\end{equation*}
The reason is simple. In the positive domain, adding something is never adding nothing. The order says that one quantity exceeds another when it differs by a strictly positive gap, and like quantities are comparable. If \(a+q=b+q\) but \(a\neq b\), one side must exceed the other by a positive gap. Adding the same \(q\) to both would then make a sum strictly exceed itself, which positivity forbids. The gap, when it exists, is unique. The theory therefore has the restricted subtraction it needs: not by introducing inverse elements or signed completion, but by reading subtraction as the unique positive gap exhibited by the order.

The result is a root discipline: the theory owns positive admitted content. Zero is absence, and negative signs are relational descriptors rather than primitive objects.

\Needspace{5\baselineskip}
\subsection{No External Background}
An important consequence of the existence axiom is that there can be no independently existing background against which reality is measured.

% LAYOUT_ONLY_GUARD:FIRST_COUNT_BACKGROUND
\Needspace{8\baselineskip}
If the primitive content of the theory is a single admitted distinction, the first count

\begin{equation*}
\begin{adjustbox}{max width=\linewidth,center}
$\displaystyle
I = 1 ,
$
\end{adjustbox}
\end{equation*}
then that admitted distinction is all that is available at the fundamental level. Nothing external to it can be assumed. An independently existing coordinate grid, background space, or absolute reference frame cannot be introduced as an additional primitive, because doing so would add structure beyond the original axiom.

Consequently, geometry, measurement, distance, orientation, and reference may enter only as later relational constructions that preserve the primitive admission discipline. Nothing may be borrowed from an external framework.

This is stronger than saying that the theory merely chooses not to assume a background coordinate system. Within this framework, such a background cannot be introduced as a theory-native primitive: any coordinate system or measurement convention must be introduced as an admitted relational construction from the primitive admission of \(I=1\).

In other words, if one distinction is all that fundamentally exists, then later mathematical and physical structures must earn their place through that relational foundation. Nothing can be assumed to exist outside it.

\Needspace{5\baselineskip}
\subsection{Measurement as Relational Comparison}
The same relational discipline applies to measurement.

The numerical values that appear in physics - integers, decimals, coefficients, and constants - are often treated as intrinsic properties of physical systems. In this framework they are read differently. A measured value is a relational readout: a comparison between admitted content and an admitted reference standard.

A unit of measurement is therefore not fundamental. It is a convention that selects a comparison standard. A gram, for example, is not a primitive physical object. When a report says that an object has a mass of ten grams, the report is comparing that object to a chosen standard about ten times over. The reference standard serves as the counting unit within that measurement convention.

The theory states this before any observer or apparatus has been constructed inside it. Which relation functions as the observer, readout locus, registration relation, or clockhood standard is built later, especially in Section 8. Section 1 fixes only the prior admissibility condition any later measurement report must satisfy. Let \(\operatorname{Adm}(x)\) mean that \(x\) is admitted as a distinguishable relation, let \(L\) denote a future readout context, let \(s\) denote the comparison standard, and let \(\operatorname{Cmp}_L(x,s)\) denote the licensed comparison relation. Then a measurement report is well defined exactly when its full comparison structure is present:

\begin{equation*}
\begin{adjustbox}{max width=\linewidth,center}
$\displaystyle
\operatorname{Meas}_L(x;s)\downarrow
\Longleftrightarrow
\operatorname{Adm}(x)
\wedge
\operatorname{Std}_L(s)
\wedge
\operatorname{Cmp}_L(x,s).
$
\end{adjustbox}
\end{equation*}
The standard itself must be admitted:

\begin{equation*}
\begin{adjustbox}{max width=\linewidth,center}
$\displaystyle
\operatorname{Std}_L(s) \Longrightarrow \operatorname{Adm}(s).
$
\end{adjustbox}
\end{equation*}
The excluded alternative is the primitive scalar: a measured value of a thing alone, belonging to no comparison.

\begin{equation*}
\begin{adjustbox}{max width=\linewidth,center}
$\displaystyle
\neg \exists q\,
\left(
\operatorname{PrimitiveScalar}(q)
\wedge
q=\operatorname{Meas}_L(x)
\right).
$
\end{adjustbox}
\end{equation*}
% LAYOUT_ONLY_GUARD:UNIT_DEFINITION
\Needspace{8\baselineskip}
A reading taken against no standard is not an imprecise measurement. It is not a measurement at all. A unit is a selected admitted standard,

\begin{equation*}
\begin{adjustbox}{max width=\linewidth,center}
$\displaystyle
\operatorname{Unit}_L(s) :=
\operatorname{Adm}(s)\wedge \operatorname{Std}_L(s),
$
\end{adjustbox}
\end{equation*}
and a reported value has the form

\begin{equation*}
\begin{adjustbox}{max width=\linewidth,center}
$\displaystyle
\operatorname{Report}_L(x;s)=q_s(x) ,
$
\end{adjustbox}
\end{equation*}
defined only when the object, the unit role, and the comparison relation are all present. This separates fundamental existence from measurement conventions. SI units and other human-scale units can communicate observer-scale reports consistently, but they are not fundamental ingredients of the internal theory. They enter only when the theory is translated into observer-scale quantities for comparison, especially in Sections 8 and 14.

\Needspace{5\baselineskip}
\subsection{Why Geometry}
The relational foundation also motivates the mathematical language used later.

Geometry is natural for a relational theory because geometry is about relations: position, orientation, adjacency, separation, order, and comparison. A point alone says almost nothing. It becomes meaningful only inside a structure that gives it relations. Geometry is therefore licensed here as relational mathematics, not as primitive ontology. The paper may use geometric language when that language preserves the admitted relations it claims to represent.

This distinction matters. The license does not say that a geometric object is primitive reality, and it does not introduce a lattice, continuum, coordinate system, metric, or physical realization at Section 1. It says that when later sections use geometry, they must use it as a faithful relational presentation rather than as an external background imposed on the theory.

The bookkeeping analogy is useful but only as an analogy. In double-entry accounting, no entry is meaningful in isolation; each posting exists through its relationship to the rest of the books. The present theory has a similar discipline, but it does not import signed ledger arithmetic. Its internal form is strictly positive transfer read through oriented sides. Relationship is the constraint; isolated quantity is not the primitive.

Maintaining that discipline keeps later mathematics from smoothing away the discrete counting basis from which the theory begins. A formalism can be internally consistent and still fail to represent the physical structure if it imports a background or treats a reporting convention as primitive content.

\Needspace{5\baselineskip}
\subsection{The Overall Program}
The program of this paper is to begin with the minimal existence axiom and develop, through forced relational constructions, the geometry later read by an embedded observer.

Many mathematical objects used in physics - negative values, fractional values, irrational quantities, imaginary quantities, and other abstractions - are not taken here as primitive features of the ontology. They may enter as mathematical language, relational structure, comparison structure, or observer-scale report, but not as starting assumptions. The starting operation is simpler:

\textbf{The starting operation is to count admitted distinctions.}

Every later scientific object must be constructed from or typed against that relational ground. This requirement does not force every exact downstream representation to be an integer.

As the theory develops, ledger closure becomes a persistence condition, and later forced relational geometry selects a definite FCC/RD realization. Observer readout and bridge comparison enter only after their own obligations have been stated and checked. The closing sections compare the count structure those later constructions force against measured physics under explicit comparison criteria.

Because the paper ends at such comparisons, its construction discipline is part of its physical content. Nothing in what follows is derived from structure introduced later than it. A numerical success found downstream is never cited as the reason an earlier rule holds. Junction, readout, bridge, and prediction constructions enter only when their obligations have been stated and checked. Naming an obligation is not discharging it.

This foundation can also lose, and it says how. Exhibit an admitted existent with no relational distinguishability - a primitive bare existence apart from any distinction - and Axiom 1 falls. Exhibit a physical measurement that is genuinely a primitive scalar, reducible to no comparison against an admitted standard, and the measurement rule falls. Exhibit a quantity domain with the stated positive-domain properties in which cancellation fails on a concrete triple, and the cancellation theorem fails. Each failure mode is specific and checkable.

The purpose of Section 1 is simply to establish this foundational ontology and mathematical discipline. Everything derived in the remainder of the paper stands on it.

